\section{Study design and data provenance}\label{sec:design}

\subsection{Need, scientific problem, and method}
The operational need is consistent long-term monitoring of
\emph{S.~alterniflora} invasion, eradication, and post-eradication
recurrence along the China coast. The underlying scientific problem is
that a single ecological state is observed through sensors whose
geometry, resolution, quality metadata, and acquisition conditions
change across decades; without an explicit contract over what
constitutes an observation, cross-era differences are confounded with
ecological change. Spartina Earth addresses this contract before any
modeling: define a stable spatial hierarchy, construct sensor-specific
events with transparent eligibility, attach tiered labels and rights,
and record provenance end to end. The present national
implementation is the mainland China coastal frame v0 plus a v1
domain-membership candidate: the administrative geometry excludes
Taiwan and the audited national label products are mainland-China
products, and this scope is stated wherever national assets appear.
The 10\,km grid (3{,}319 fixed cells) is deliberately distinct from
domain membership (Section~\ref{sec:national}: 3{,}011 active
mainland-coastal members, 298 excluded artifacts, ten provisional);
neither is frozen.

\subsection{Staged design and authority boundaries}
The program proceeds in stages with explicit stop gates (Figure~\ref{fig:arch}).
\emph{M2.1a} established region-of-interest provenance and a
metadata-only census; \emph{M2.1a2-R1} corrected the observation-unit
and quality-attribution defects documented here; \emph{M2.1b} (Issue
\#13), now executed, was a deliberately small, eight-cell real-pixel
pilot; national
frame construction (Issue \#14) proceeds as metadata-only work; the
full Zhejiang dataset, benchmark freeze, GOLD validation, and any
representation-learning experiments are later stages. Nothing in this
paper assumes the output of a later stage.

\begin{figure}[t]\centering
\figplaceholder{Figure 1}{Spartina Earth hierarchy:
China coast $\rightarrow$ provisional bay envelope $\rightarrow$ fixed
10\,km analysis cell (cell\_geometry vs bay\_clip\_geometry)
$\rightarrow$ observation event / acquisition group $\rightarrow$
optical, SAR and context modalities $\rightarrow$ label tier, time,
tide, rights, provenance. Framework figure; no metric values.}
\label{fig:arch}
\end{figure}

\subsection{Provenance chain and reproducibility}
Every derived product is traceable as
\texttt{source\ scene\ $\rightarrow$ preprocessing $\rightarrow$ label
$\rightarrow$ model/config $\rightarrow$ checkpoint $\rightarrow$
inference\ config $\rightarrow$ output}, with a recorded Git commit,
configuration hash, data manifest version, and seed for any executed
run. Manifests (CSV/Parquet, row-level fingerprints) are tracked in
Git; raster bytes, arrays, and checkpoints live outside Git.
Superseded artifacts are never overwritten: a new version is written
alongside the old file, which is retained with a
\texttt{.SUPERSEDED.json} sidecar stating the reason. Numbers from
legacy papers, spreadsheets, or code are treated as unverified until
reproduced from raw evidence; if an item cannot be found it is recorded
as \textsc{unknown}/\textsc{missing}, never inferred.

\subsection{Spatial leakage control}
Train/validation/test membership is defined at the cell level on a
fixed metric grid, never by random patch draws. Scenes of one location
across dates stay grouped; candidate held-out tracks (cross-bay
transfer, held-out era, sensor-generation shift, missing-modality) are
proposed from census evidence before any model is trained. The
national extension adds target-independent coastal relevance so that
label history never defines who enters the sampling frame.
